% Stage C2: the bit arrays of 64-bit words, the pieces, the masks, the sentinel scan and the walk by
% trailing zeros, as a data refinement of the plain sieve.
\chapter{The bit arrays of words}\label{chap:words}

Stage~C2 replaces the Booleans of the plain sieve by the bit arrays of 64-bit words of the C code:
the numerators of a denominator are cut into pieces of at most \verb|array_size| bit arrays, the
rows of the moduli are and-ed into the words of a piece, masks clear the bits beyond the two ends
of the numerator range, a scan steps over the empty bit arrays up to a sentinel, and the set bits
of the others are visited through their trailing zeros. This is a data refinement in the sense of
Chapter~\ref{chap:sieve}: the sieve of a denominator returns what the brute-force reference finds
over the numerators that survive the moduli, now computed with words, so the equation of
Theorem~\ref{thm:findPoints_eq_bruteForce} and the main theorems hold with their statements
unchanged.

The C code has one bit-array type per build. The model takes the width of the default build, four
words per bit array, as a constant that no proof depends on, so that the proofs hold for every
width; and it takes the first stage of the builds that and the rows into one bit array at a time,
which the default build does in chunks of sixteen registers. So the model is the build that
ratpoints tests as \verb|chunk1| (256-bit bit arrays and \verb|-DRATPOINTS_CHUNK=1|, one of the
configurations of \verb|make test4configs|). The numerators keep the plain layout of stage~C1, one
bit per numerator from $-H$ to $H$; the packing by the 2-adic classes comes with stages~E and~F, the
chunks and the second and third stages of the sieve with stage~F.

\section{What the C code does with words}\label{sec:wordsC}

A bit array has \verb|RBA_LENGTH| bits in \verb|RBA_PACK| words of 64 bits, and the build fixes
its type and its operations (\csource{rp-private.h}{113}{300}): \verb|AND|, the tests
\verb|TEST| and \verb|TESTZ| for a nonzero and a zero bit array, the masks \verb|MASKL| and
\verb|MASKU| that clear the lowest and the highest $s$ bits for $0 \le s <$ \verb|RBA_LENGTH|, and
the reads \verb|EXT0| and \verb|EXT| of a word; bit $t$ of a bit array is bit $t \bmod 64$ of its
word $\lfloor t/64 \rfloor$. The default build has 256-bit bit arrays, four words each
(\csource{rp-private.h}{167}{224}); the build with plain 64-bit words has one. Write $R$ for
\verb|RBA_LENGTH|. The C code calls the index of a bit array its \emph{word number}. In the plain
layout bit $t$ of the bit array with word number $w$ stands for the numerator $R w + t$; ratpoints
uses this layout for the classes of denominators with stride $1$, and otherwise bit $t$ stands for
$a_0 + 2^k(R w + t)$, for the stride $2^k$ and the offset $a_0$ of the 2-adic class of $b$. For a
denominator $b$, \verb|sift| (\csource{find_points.c}{3442}{3617}):
\begin{enumerate}
  \item sets up, once per denominator, at the first interval that \verb|interval_bits| does not
    report empty, the row of $b\,2^{-k} \bmod p$ (the row of $b$ in the plain layout) in the table
    of each modulus of the first two stages, building the rows that no earlier denominator needed
    (\csource{find_points.c}{3490}{3561}, with \verb|set_rows|,
    \csource{find_points.c}{3149}{3170}): the row of a residue $r$ modulo $p$ is a table of $p$ bit
    arrays, the one for word number $w$ at index $(w + o) \bmod p$, where the offset $o$ is a
    multiple of $p$ in the plain layout, and carries $a_0 (2^k R)^{-1} \bmod p$ otherwise
    (\csource{find_points.c}{1142}{1183});
  \item for each interval $[l, u]$ of the search region (the intervals that the Sturm step keeps
    within the domain, which is $[-h, h]$ with $h \ge H$ when the caller gives none:
    \csource{find_points.c}{4829}{4892}; stage~H), takes the numerators
    $\lceil b l \rceil, \dots, \lfloor b u \rfloor$ within $[-H, H]$ as the bits
    $[\mathit{low}, \mathit{high}]$ of the packing of the class of $b$ (\verb|interval_bits|,
    \csource{find_points.c}{3410}{3434}; \verb|numerators_to_bits|,
    \csource{find_points.c}{3342}{3370}), returning at the first interval beyond $H$ and skipping
    an empty one (\csource{find_points.c}{3471}{3488}, \csource{find_points.c}{3568}{3568}); turns
    $[\mathit{low}, \mathit{high}]$ into the word numbers
    $w_{\mathrm{low}} = \lfloor \mathit{low}/R \rfloor$ and
    $w_{\mathrm{high}} = \lfloor \mathit{high}/R \rfloor + 1$, and cuts these into pieces of at most
    \verb|array_size| bit arrays, the first with the mask $\mathit{low} \bmod R$ on its first bit
    array and the last with the mask $R - 1 - (\mathit{high} \bmod R)$ on its last one
    (\csource{find_points.c}{3563}{3612}); for a class of stride $1$ and the search region
    $[-H, H]$ (no Sturm step, as under \verb|-s|, or one that keeps the whole domain),
    $\mathit{low} = -H$ and $\mathit{high} = H$;
  \item calls \verb|_ratpoints_sift0| (\csource{sift.c}{498}{1305}) on each piece, which and-s the
    rows of the moduli of the first stage into the bit arrays of the piece, in the builds with one
    register per chunk after writing the 2-adic pattern of the class into them
    (\csource{sift.c}{943}{1022}), and in the default build in chunks of sixteen registers, the
    first modulus anding the pattern in, with legs over the tail of the piece
    (\csource{sift.c}{575}{942}, \csource{sift.c}{452}{496}); clears the masked bits
    (\csource{sift.c}{1026}{1041}); writes a sentinel just past the piece and scans for the
    nonzero bit arrays (\csource{sift.c}{1061}{1092}); and-s into each the rows of the moduli of
    the second stage (\csource{sift.c}{1174}{1202}); and visits the set bits of its words from the
    lowest up through their trailing zeros (\csource{sift.c}{262}{264},
    \csource{sift.c}{1213}{1299}), each numerator going through the test for common factors and
    the third stage (\csource{sift.c}{433}{446}) to the exact check
    (\csource{find_points.c}{5250}{5338}).
\end{enumerate}
\verb|array_size| is a parameter of the library; the program always passes 256
(\csource{main.c}{347}{347}, \csource{main.c}{221}{222}), and \verb|find_points| takes 256 for a
value $\le 0$, caps it at $2 \lceil H/64 \rceil$ and at $2^{20}$, and rounds it up to a multiple of
the chunk (\csource{find_points.c}{4666}{4677}).

\section{The constants and the parameters}\label{sec:wordsParams}

\begin{definition}[The width of the bit arrays]\label{def:rbaPack}
  \lean{Ratpoints.rbaPack, Ratpoints.rbaLength, Ratpoints.rbaPack_pos}
  \leanok
  A bit array has $\mathrm{rbaPack} = 4$ words of 64 bits, $\mathrm{rbaLength} = 256$ bits: the
  values of the default build (\csource{rp-private.h}{174}{175} and
  \csource{rp-private.h}{172}{173}, with the word length \csource{rp-private.h}{49}{51}); the
  builds with 512-bit, 128-bit and 64-bit bit arrays have 8, 2 and 1 words
  (\csource{rp-private.h}{127}{128}, \csource{rp-private.h}{233}{234},
  \csource{rp-private.h}{259}{260}, \csource{rp-private.h}{283}{284}; their lengths:
  \csource{rp-private.h}{125}{126}, \csource{rp-private.h}{231}{232},
  \csource{rp-private.h}{257}{258}, \csource{rp-private.h}{281}{282}), and the build is chosen by
  the lines \csource{rp-private.h}{113}{113}, \csource{rp-private.h}{167}{167},
  \csource{rp-private.h}{226}{226}, \csource{rp-private.h}{252}{252},
  \csource{rp-private.h}{278}{278}, \csource{rp-private.h}{300}{300} and
  \csource{rp-private.h}{55}{58}. The vector type of a bit array becomes $\mathrm{rbaPack}$
  consecutive words of one array of \verb|UInt64|, the length is computed as
  $64 \cdot \mathrm{rbaPack}$, and the shifts by \verb|RBA_SHIFT| and \verb|LONG_SHIFT| and the
  masks by \verb|RBA_LENGTH - 1| and \verb|LONG_MASK| have no counterpart of their own: the model
  divides by $\mathrm{rbaLength}$ and by 64 and takes remainders. The proofs use only
  $\mathrm{rbaPack} > 0$, so they hold for every build.
\end{definition}

\begin{definition}[The pieces]\label{def:arraySize}
  \uses{def:Params}
  \lean{Ratpoints.defaultArraySize, Ratpoints.maxArraySize, Ratpoints.Params.arraySize, Ratpoints.Params.arraySizeAt, Ratpoints.defaultArraySize_pos, Ratpoints.maxArraySize_pos, Ratpoints.Params.one_le_arraySizeAt}
  \leanok
  The parameters have the field $\mathrm{arraySize}$, at most how many bit arrays a piece has, with
  $0$ for the default $\mathrm{defaultArraySize} = 256$ (\csource{ratpoints.h}{354}{356}). For the
  height bound $H$ the sieve takes $\mathrm{arraySizeAt}(H)$: the field, or the default for $0$,
  capped at $2 \lceil H/64 \rceil$ (with the macro \verb|CEIL|, \csource{rp-private.h}{87}{87}) and
  at $\mathrm{maxArraySize} = 2^{20}$ (\csource{find_points.c}{124}{124}), with no rounding to a
  multiple of the chunk, which is one bit array in the build of the model
  (\csource{find_points.c}{4666}{4677}). It is at least $1$ for $H \ge 1$, which the proofs read off
  $0 < \mathrm{defaultArraySize}$ and $0 < \mathrm{maxArraySize}$, two facts about the constants
  checked once.
\end{definition}

The parameters of the executable take the default, as the program does: \verb|sieveParams|
(Definition~\ref{def:sieveParams}) gives the field the value $\mathrm{defaultArraySize}$, as the C
lines of the program that set it (\csource{main.c}{347}{347}, \csource{main.c}{221}{222}). The field
is a free choice of the library that changes neither the points found nor their order, so the
theorems quantify over it as over the other parameters, and \verb|Params.WF| does not constrain it:
every value gives pieces of at least one bit array. The comparator challenge imports \verb|Params|,
so its structure has this field, while the statements of the challenge and \verb|Params.WF| are
those of stage~C1. The number of registers of a chunk, \verb|RATPOINTS_CHUNK|, becomes a constant in
stage~F, together with the chunked first stage and the copies in the rows that it decides; for the
model it is $1$.

\section{The rows as words}\label{sec:wordRows}

\begin{definition}[The row of a denominator in words]\label{def:row}
  \uses{def:admits, def:rbaPack}
  \lean{Ratpoints.Sieve.packedPattern, Ratpoints.Sieve.window, Ratpoints.Sieve.Entry.row, Ratpoints.Sieve.Entry.admits_congr, Ratpoints.Sieve.Entry.admits_add_mul, Ratpoints.Sieve.getLsbD_window, Ratpoints.Sieve.Entry.getLsbD_row}
  \leanok
  For the table of a prime $p$ and a denominator $b$: the pattern of one period, whether the row of
  $b$ admits $N$ for $N = 0, \dots, p - 1$ (Definition~\ref{def:admits}); the pattern over the
  positions $0, \dots, p + 63$, packed into words; for $j < p$, word $j$ of the row, the 64 bits of
  the packed pattern from position $64 j \bmod p$ on, cut out with two shifts; and these words
  repeated $\mathrm{rbaPack}$ times, which makes $p$ bit arrays. Since the row admits $a$ exactly
  when it admits $a + p k$, bit $t$ of word $j$ says whether the row admits $64 j + t$, for every
  $j < p \cdot \mathrm{rbaPack}$. ratpoints builds the same table from the table of the prime and
  the inverses, the pattern of one period bit by bit and then the words by shifting it through
  them, by a function for each prime that its build generates (\csource{init.c}{42}{46},
  \csource{init.c}{62}{83}, \csource{init.c}{118}{246}, \csource{gen_init_sieve_h.c}{48}{62}), with
  the copies \csource{init.c}{99}{111}, and takes the row of $0$ from the generated pattern
  \verb|sieves0| (\csource{gen_find_points_h.c}{102}{134}); these ranges are shared with
  Definition~\ref{def:admits} until the faithful builders of stage~F. The copies of the first bit
  arrays at the end of a row, which only the chunked first stage reads, are none in the build of
  the model.
\end{definition}

\begin{definition}[The rows of a modulus]\label{def:rows}
  \uses{def:row}
  \lean{Ratpoints.Sieve.Entry.rows, Ratpoints.Sieve.Entry.get_rows, Ratpoints.Sieve.Entry.row_mod, Ratpoints.Sieve.SieveSpec}
  \leanok
  The rows of the residues $r = 0, \dots, p - 1$, each made when a denominator first needs it, as a
  suspended computation that runs once, and kept for the run, which makes these computations once
  (Definition~\ref{def:sieveRun}). They stand for the field \verb|sieve| of the sieve entry of
  ratpoints, whose row of $0$ \verb|examine_prime| installs from the build
  (Definition~\ref{def:examinePrime}) and whose other rows \verb|sift| builds when first needed; the
  model makes the row of $0$ when first needed too. The row of $b$ is that of $b \bmod p$. For the
  denominator at hand, each modulus comes with its row, as in the \verb|sieve_spec| of the C code
  (\csource{rp-private.h}{378}{379}), without the offset and the pointers into the row.
\end{definition}

\section{Words and bit arrays}\label{sec:wordOps}

The words are \verb|UInt64|, and the proofs read their bits through \verb|BitVec.getLsbD| of
\verb|UInt64.toBitVec|, whose lemmas for the bitwise operations Lean core has, so that no
\verb|bv_decide| is needed. The bit arrays of a piece and of a row lie in one array of words, bit
array $j$ in the words $\mathrm{rbaPack} \cdot j, \dots, \mathrm{rbaPack} \cdot (j + 1) - 1$; the
macros of the C code are the word operations they implement, each one definition for every
build.

The loops of the C code that end on the data are Lean's \verb|while| and \verb|repeat|, as in the C
code: the lowest set bit, the walk through the set bits, the scan up to the sentinel and the loop
over the pieces. Their logical meaning is a fixed point of the loop body, so their contracts state
total correctness and prove that each loop ends, by a measure: $64$ less the count for the lowest
set bit (the count stays at most the trailing zeros of the nonzero word), the word itself for the
walk (clearing its lowest set bit makes it smaller), the distance to the sentinel for the scan, and
the distance to $w_{\mathrm{high}}$ for the pieces. Hence the hypothesis $A \ge 1$, only for the
loop over the pieces, and the fact $\mathrm{rbaPack} > 0$, checked once, which makes the sentinel
nonzero and gives the bit arrays their numerators.

\begin{definition}[Operations on bit arrays]\label{def:bitArrayOps}
  \uses{def:rbaPack}
  \lean{Ratpoints.Sieve.andBitArray, Ratpoints.Sieve.clearLowBits, Ratpoints.Sieve.clearHighBits, Ratpoints.Sieve.isZeroAt, Ratpoints.Sieve.getElem!_andBitArray, Ratpoints.Sieve.getLsbD_clearLowBits, Ratpoints.Sieve.getLsbD_clearHighBits, Ratpoints.Sieve.isZeroAt_eq_true_iff}
  \leanok
  Bit array $j$ of an array of words and-ed with bit array $x$ of a row, word by word (\verb|AND|:
  \csource{rp-private.h}{176}{176}, and \csource{rp-private.h}{134}{134},
  \csource{rp-private.h}{237}{237}, \csource{rp-private.h}{263}{263},
  \csource{rp-private.h}{287}{287}); the lowest $s$ bits of bit array $0$ cleared, by zeroing the
  words below $\lfloor s/64 \rfloor$ and masking the next one (\verb|MASKL|:
  \csource{rp-private.h}{206}{213}, and \csource{rp-private.h}{153}{155},
  \csource{rp-private.h}{242}{244}, \csource{rp-private.h}{268}{270},
  \csource{rp-private.h}{291}{291}); the highest $s$ bits of bit array $j$ cleared, in the same way
  from the top (\verb|MASKU|: \csource{rp-private.h}{214}{221}, and
  \csource{rp-private.h}{156}{159}, \csource{rp-private.h}{245}{247},
  \csource{rp-private.h}{271}{273}, \csource{rp-private.h}{292}{292}); and whether bit array $j$ is
  zero (\verb|TEST|, \verb|TESTZ|: \csource{rp-private.h}{179}{202}, and
  \csource{rp-private.h}{137}{152}, \csource{rp-private.h}{241}{241},
  \csource{rp-private.h}{267}{267}, \csource{rp-private.h}{290}{290},
  \csource{rp-private.h}{306}{308}). The low mask keeps bit $t$ of word $y$ exactly when
  $s \le 64 y + t$, and for $s < R$ the high mask on bit array $j$ keeps bit $t$ of word
  $y < \mathrm{rbaPack} \cdot (j + 1)$ exactly when $64 y + t + s < R (j + 1)$. The masks follow the
  form of the 512-bit build, of which the macros of the other builds are unrolled cases. Of the
  other lines of the blocks of the builds, the constant \verb|zero| and the reads \verb|EXT0| and
  \verb|EXT| go with the scan (Definition~\ref{def:scanPiece}); the constructor \verb|RBA|, used
  only for the 2-adic pattern of a class, waits for stage~E; the includes of the vector intrinsics
  and the printing of bit arrays are excluded from the model.
\end{definition}

\begin{definition}[The lowest set bit]\label{def:ctz}
  \lean{Ratpoints.Sieve.ctz, Ratpoints.Sieve.ctz_spec, UInt64.getLsbD_and_sub_one, UInt64.toNat_and_sub_one_lt}
  \leanok
  The position of the lowest set bit of a nonzero word, by the loop of \verb|RP_CTZL| for a
  compiler without the builtin: shift the word right while its lowest bit is clear, counting the
  steps (\csource{rp-private.h}{429}{444}). Its contract: for $x \ne 0$ the result is the number
  of trailing zeros of $x$. Clearing the lowest set bit, $x \mathbin{\&} (x - 1)$, keeps every
  other bit of $x$, and for $x \ne 0$ it makes $x$ smaller.
\end{definition}

\section{One piece}\label{sec:piece}

Fix a piece: the bit arrays with the word numbers $w_0, \dots, w_0 + n - 1$, in an array of
$(n + 1) \cdot \mathrm{rbaPack}$ words, the last bit array for the sentinel. Word $y$ of the
piece, for $y < n \cdot \mathrm{rbaPack}$, holds the numerators $R w_0 + 64 y + t$ for
$t < 64$; this is what makes the width disappear from the statements. The components of a piece
receive the moduli with their rows for $b$ and the products of the check rather than the tables
and the coefficients, so their contracts quantify over the tables whose rows they are, as the
contract of the check of stage~C1 quantifies over the test of the numerators.

\begin{definition}[The first stage]\label{def:andRows}
  \uses{def:bitArrayOps, def:rows}
  \lean{Ratpoints.Sieve.andRows, Ratpoints.Sieve.andRows_spec}
  \leanok
  Every word of the piece set to all ones, where ratpoints writes the 2-adic pattern of the class of
  $b$ (stage~E); then for each modulus $p$ with its row, for $r = (-w_0) \bmod p$, bit array $j$ is
  and-ed with bit array $(w_0 + j) \bmod p$ of the row, in the three loops of the C code: $n$ bit
  arrays of the row from $p - r$ on when $n < r$, and otherwise $r$ bit arrays from $p - r$ on, then
  whole periods of $p$ from $0$, then the rest from $0$ (\csource{sift.c}{943}{1022}). Every modulus
  is applied here, those of the second and the third stage too (\csource{sift.c}{291}{312},
  \csource{sift.c}{1083}{1086}, \csource{sift.c}{1174}{1202}). Its contract: for every list of
  tables with $p > 0$ whose rows for $b$ the moduli carry, bit $t$ of word
  $y < n \cdot \mathrm{rbaPack}$ is $t < 64$ and the survivor test of $R w_0 + 64 y + t$.
\end{definition}

\begin{definition}[The check of a numerator]\label{def:checkNum}
  \uses{def:bruteForce, def:swapPoint}
  \lean{Ratpoints.Sieve.checkNum, Ratpoints.Sieve.pointsOfDen_eq_flatMap_checkNum, Ratpoints.Sieve.checkNum_eq_map, Ratpoints.Sieve.flatMap_checkNum}
  \leanok
  The test for common factors of $a$ and $b$ and the exact check, from the products $c_k b^{d - k}$
  and the factor $b$ for odd $d$, which come in as arguments: the points $(a, y, b)$ and, for
  $y \ne 0$, $(a, -y, b)$, each mapped by $\mathrm{swapPoint}$ (Definition~\ref{def:swapPoint}) when
  the flag of the reversal, an argument before the numerator, is set (\csource{sift.c}{202}{205},
  \csource{sift.c}{318}{343}, \csource{sift.c}{438}{446}, \csource{find_points.c}{5250}{5338}). So
  the check with the flag set gives the points of the check with the flag clear, mapped, and so do
  the checks of a list of numerators, one after the other. What the brute-force reference finds over
  $b$ and a list of numerators is the concatenation of the checks with the flag clear over the list,
  with the products and the factor of $c$, $d$ and $b$.
\end{definition}

\begin{definition}[The set bits of a word]\label{def:walkWord}
  \uses{def:ctz, def:checkNum}
  \lean{Ratpoints.Sieve.walkWord, Ratpoints.Sieve.walkWord_spec}
  \leanok
  For a word $x$ whose bit $t$ stands for the numerator $a_0 + t$: while $x \ne 0$, the check of
  $a_0 + \mathrm{ctz}(x)$, with the flag of the reversal that the walk hands on, and
  $x \mathbin{\&} (x - 1)$ for $x$: the loop of the C code with its body at both places of the
  extraction (\csource{sift.c}{262}{264}, \csource{sift.c}{1227}{1255},
  \csource{sift.c}{1265}{1292}). Its contract: the checks of the numerators $a_0 + t$ for the set
  bits $t$ of $x$, in increasing order, are appended.
\end{definition}

\begin{definition}[The scan and the extraction]\label{def:scanPiece}
  \uses{def:walkWord, def:bitArrayOps}
  \lean{Ratpoints.Sieve.scanPiece, Ratpoints.Sieve.scanPiece_spec}
  \leanok
  A sentinel, a bit array of ones, just past the piece; then from the first bit array on: step over
  the zero bit arrays, stop at the sentinel, and walk the words of each other bit array, with the
  flag of the reversal, the bit $0$ of word $\kappa$ of bit array $i$ standing for the numerator
  $R (w_0 + i) + 64 \kappa$ (\csource{sift.c}{1061}{1082}, \csource{sift.c}{1088}{1092},
  \csource{sift.c}{1213}{1226}, \csource{sift.c}{1257}{1264}, \csource{sift.c}{1293}{1299}; with
  \verb|zero|, \csource{rp-private.h}{133}{133}, \csource{rp-private.h}{205}{205},
  \csource{rp-private.h}{236}{236}, \csource{rp-private.h}{262}{262},
  \csource{rp-private.h}{286}{286}, and \verb|EXT0|, \verb|EXT|, \csource{rp-private.h}{135}{136},
  \csource{rp-private.h}{177}{178}, \csource{rp-private.h}{238}{239},
  \csource{rp-private.h}{264}{266}, \csource{rp-private.h}{288}{289}). Its contract: if bit $t$ of
  word $y$ is $t < 64$ and $k(R w_0 + 64 y + t)$, for a test $k$ of the numerators and every
  $y < n \cdot \mathrm{rbaPack}$, the result is the concatenation of the checks of the numerators
  $R w_0, \dots, R(w_0 + n) - 1$ that pass $k$, in increasing order.
\end{definition}

\begin{definition}[One piece]\label{def:siftPiece}
  \uses{def:andRows, def:scanPiece}
  \lean{Ratpoints.Sieve.siftPiece, Ratpoints.Sieve.siftPiece_spec}
  \leanok
  A fresh array of $n + 1$ bit arrays, as the worker's array of \verb|array_size| bit arrays and two
  more (\csource{find_points.c}{3119}{3121}); the first stage; the low mask $m_l$ on the first bit
  array if $m_l \ne 0$ and the high mask $m_h$ on the last if $m_h \ne 0$
  (\csource{sift.c}{1030}{1041}); and the scan, with the flag of the reversal. The piece also has
  the head of \verb|_ratpoints_sift0| (\csource{sift.c}{498}{531}), the count of its words
  (\csource{sift.c}{548}{548}), the timer and the test before its second stage
  (\csource{sift.c}{1056}{1058}) and its end (\csource{sift.c}{1300}{1305}). Its contract, for
  $m_h < R$ and every list of tables with $p > 0$ whose rows for $b$ the moduli carry: the
  concatenation of the checks of the numerators from $R w_0 + m_l$ to $R (w_0 + n) - 1 - m_h$ that
  survive the tables.
\end{definition}

\section{One denominator and the run}\label{sec:siftDen}

\begin{definition}[One denominator]\label{def:sift}
  \uses{def:siftPiece}
  \lean{Ratpoints.Sieve.sift, Ratpoints.Sieve.sift_spec}
  \leanok
  The row of $b$ for each modulus, for the rows of the first two stages
  (\csource{find_points.c}{3149}{3170}, \csource{find_points.c}{3490}{3561}) and the table of the
  third (\csource{sift.c}{405}{431}); the products $c_k b^{d - k}$ and the factor of the check,
  which also gets the flag of the reversal; the word numbers
  $w_{\mathrm{low}} = \lfloor -H/R \rfloor$ and $w_{\mathrm{high}} = \lfloor H/R \rfloor + 1$; and
  while $w_0 < w_{\mathrm{high}}$, the piece from $w_0$ to $w_1 = \min(w_0 + A, w_{\mathrm{high}})$,
  with the low mask $-H - R w_{\mathrm{low}}$ if it is the first and the high mask
  $R - 1 - (H \bmod R)$ if it is the last (\csource{find_points.c}{3442}{3488},
  \csource{find_points.c}{3563}{3617}). Its contract, for $A \ge 1$ and the rows of tables with
  $p > 0$: what the brute-force reference finds over $b$ and the numerators $-H, \dots, H$ that
  survive the tables, the result of the sieve of $b$ followed by its check in stage~C1, mapped by
  $\mathrm{swapPoint}$ when the flag is set. The masked ranges of the pieces lie end to end and make
  up $[-H, H]$: the low mask moves the start of the first piece from $R w_{\mathrm{low}}$ to $-H$,
  the high mask the end of the last one from $R w_{\mathrm{high}} - 1$ to $H$, and a piece between
  them has no numerator outside $[-H, H]$.
\end{definition}

\begin{lemma}[The words compute the Booleans]\label{lem:sift_eq_checkDen}
  \uses{def:sift, def:sieveDen, def:checkDen}
  \lean{Ratpoints.Sieve.toList_run_sift_eq_checkDen}
  \leanok
  For $A \ge 1$ and tables with $p > 0$, the points of the denominator $b$ that \verb|sift| returns
  are those that the check of stage~C1 returns on the sieve of stage~C1, mapped by
  $\mathrm{swapPoint}$ when the flag of the reversal is set.
\end{lemma}

\begin{proof}
  \leanok
  Both are what the brute-force reference finds over $b$ and the numerators $-H, \dots, H$ that
  survive the tables, the points of \verb|sift| mapped when the flag is set, by the contract of
  Definition~\ref{def:sift} and those of Definitions~\ref{def:sieveDen} and~\ref{def:checkDen}.
\end{proof}

The sieve of one denominator and the check of stage~C1 are the reference of the words, with this
lemma, and hold no C ranges: the components of this chapter model those lines, with the cuts above.
The run over the denominators (Definition~\ref{def:sieveRun}) makes the rows of the moduli once and
returns, for each denominator, the points that \verb|sift| finds; its contract holds for $A \ge 1$
and tables with $p > 0$, and \verb|sift| computes the residues $b \bmod p$ once per denominator,
where ratpoints computes the residues of $b / 2^k$ before each call of \verb|sift|
(\verb|fill_bp_list|). \verb|findPoints| (Definition~\ref{def:findPoints}) runs it with
$A = \mathrm{arraySizeAt}(H)$ of its parameters, the tables of its moduli, which are primes, and the
flag of the reversal clear; so the proof of Theorem~\ref{thm:findPoints_eq_bruteForce} needs two
more facts, the lower bound of the array size and the positivity of the moduli, and the statements
of the theorems and of the comparator challenge are those of stage~C1.

Each component of this chapter differs from the C code in a stated way, recorded in its
\verb|Deviation:| paragraph:
\begin{itemize}
  \item the plain layout, bit $t$ of the bit array with word number $w$ for the numerator
    $R w + t$, every numerator from $-H$ to $H$ and no 2-adic pattern, for the packing of the
    2-adic classes (E, F) and the search region (H);
  \item every modulus in the first stage, over every bit array, for the second stage, which and-s
    the rows of its moduli, one after the other, into each bit array that the scan finds nonzero,
    and stops as soon as no bit of it is left, and the third stage, which tests one numerator at a
    time after the test for common factors, stops at the first prime that rejects it, skips the
    primes that divide $b$, and looks up the inverses of $b$ when the first numerator of $b$ reaches
    it (F);
  \item the first stage of the builds with one register per chunk, for the chunks of sixteen
    registers of the default build with the legs over the tail of a piece (F);
  \item the rows from Definition~\ref{def:admits}, one period laid out in words, for the builders
    of the C code with their bit arithmetic (F);
  \item the rows of all residues, that of $0$ among them, made when a denominator first needs them
    and kept for the run, where ratpoints installs the row of $0$ from the build when it examines
    the prime and keeps the rows in its sieve entry (F);
  \item the row of $b \bmod p$, computed in \verb|sift|, for the row of $b\,2^{-k} \bmod p$ that
    \verb|set_rows| reads from the residues that \verb|fill_bp_list| computes before the call, and
    the row of the word number $w$ at index $w \bmod p$, where ratpoints adds the offset of the
    packing and a multiple of $p$ that keeps the index non-negative, and reduces by a
    multiplication when that is exact (F);
  \item indices for the pointers of the C code, into the rows and into the survivors, and a scan
    that holds its position as an index and reads the words of a bit array in place, where ratpoints
    recovers the word number from a pointer difference and copies the bit array;
  \item the lowest set bit by the loop of the C code for a compiler without the builtin;
  \item a fresh array for each piece, for the survivors array of the worker that every piece reuses,
    with one bit array more for the alignment, and for each row, for the space that the builders
    take from the buffer of the tables;
  \item a plain $\gcd$ for the binary $\gcd$ of \verb|relprime|, whose loop body has no branch (F);
  \item the products $c_k b^{d - k}$ computed in \verb|sift| for every denominator and passed to
    the check, where ratpoints computes them in the check when the first survivor of $b$ reaches
    it;
  \item a list of points for the callback \verb|process| and for the counts of points that the C
    functions return, with no stop it may request, no \verb|NO_Y| and no \verb|NO_CHECK| (G), and
    the flag of the reversal as an argument of the check, which the components from the run down to
    it hand on, for the bit \verb|RATPOINTS_REVERSED| of the flags that the check reads;
  \item no counts of the denominators sieved, of the words swept and of the bit arrays, bits and
    numerators that reach each stage (F), no threads, and natural numbers and integers for the
    \verb|long|s, with a \verb|Nat| for \verb|array_size| so that only $0$ stands for the default
    (the range obligations: D, F).
\end{itemize}

\section{The command-line program and the differential tests}\label{sec:wordsCli}

\verb|ratpoints-lean| runs \verb|findPoints| as before, now with the array size 256 of the
program, capped at $2 \lceil H/64 \rceil$ (Definition~\ref{def:arraySize}). At four words per bit
array every denominator up to $H = 32767$ is then a single piece; a denominator with several
pieces, and masks on different pieces, occur in the runs of \verb|test4.sh| from height $40000$ on
and in the \verb|#guard| examples with small array sizes. Compiled, with the 30 default moduli on a
curve of degree~$6$, \verb|ratpoints-lean| takes $0.027$~seconds at $H = 1000$, $0.21$ at
$H = 5000$ and $6.6$ at $H = 30000$, against $0.16$ and $3.6$ for the plain sieve of stage~C1
(Section~\ref{sec:sieveCli}), which would take about $130$ at $H = 30000$; at $H = 100$ it takes
$0.015$~seconds against $0.007$, since every row that a denominator needs is made whole, $p$ bit
arrays, while a small height reads only a few of them. The C program stays below $0.01$~seconds at
these heights. So the differential tests search up to height $30000$ in every suite: the six runs of
\verb|test4.sh| at height $30000$ join the suite \verb|test4|, and the suite of point-rich curves
adds the height $30000$. The 13 runs of \verb|test4.sh| with heights from $40000$ to $100000$ form
the suite \verb|test4high|, which runs locally, not in CI, in about four minutes.

\section{Files and names}

The constants $\mathrm{rbaPack}$, $\mathrm{rbaLength}$, $\mathrm{defaultArraySize}$ and
$\mathrm{maxArraySize}$ are in \verb|Ratpoints/Model/Constants.lean|, with the facts, checked by
\verb|decide|, that $\mathrm{rbaPack}$, $\mathrm{defaultArraySize}$ and $\mathrm{maxArraySize}$ are
positive; the field and the cap are in \verb|Ratpoints/Model/Params.lean|, and the components in
\verb|Ratpoints/Model/Sieve.lean|, in the namespace \verb|Ratpoints.Sieve| after those of stage~C1,
with the run below them, each with \verb|#guard| examples. The lower bound of the array size is in
\verb|Ratpoints/Proof/Params.lean|. The contracts, without Mathlib, are in
\verb|Ratpoints/Proof/Words.lean|, with that of the run and its equation, which
\verb|Ratpoints/Proof/FindPoints.lean| uses; the lemmas on vectors, arrays, lists, ranges, natural
numbers, integers and words that are not about the model are in \verb|Ratpoints/Proof/Bits.lean|,
under their names in the namespaces of Lean core.

\section{The evidence}\label{sec:words-evidence}

Besides the proofs:
\begin{itemize}
  \item a reading of the C code of the word sieve at the pin, with simulations of the pieces and
    their masks for four widths, both arms of the first stage, the row builders and the whole path
    of \verb|sift| and \verb|_ratpoints_sift0| against the list of survivors, one numerator at a
    time, on 10400 cases, which fixed Section~\ref{sec:wordsC}; and two independent reviews of the
    C ranges and the deviations of the components against the C code, one before the stage and one
    after it, each finding checked by a second reader, which fixed the owners of the lines of the
    builds (those of \verb|RBA_LENGTH| with $\mathrm{rbaLength}$, those of \verb|RBA_PACK| with
    $\mathrm{rbaPack}$), the cuts of the extraction between the scan and the walk, the set-up of
    the third stage in \verb|sift| and the count of the words of a piece in \verb|siftPiece|, and
    completed the deviations (the allocation and the keeping of the rows, the second stage, the
    factor $b$ for odd degree, the \verb|long|s);
  \item the \verb|#guard| examples of the model: the masks, the lowest set bit, the rows against
    the row test and against \verb|sieves0|, the first stage against the survivor test, the walk
    and the scan against the set bits of their words, the check of a numerator against the
    reference, one piece with its masks around the numerator of a point, \verb|sift| against the
    plain sieve with pieces of one, two and three bit arrays, and \verb|findPoints| against the
    brute-force reference with pieces of one, two and three bit arrays at heights around the
    multiples of 64 and of $\mathrm{rbaLength}$;
  \item the differential tests of Section~\ref{sec:wordsCli}: \verb|ratpoints-lean| and the C
    program agree on all 3155 comparisons of the eight default suites and on the 13 of
    \verb|test4high|;
  \item Comparator accepts the solution of the challenge, whose statements are those of
    stage~C1;
  \item the project builds, with all its proofs, for other values of the constants of the build: 1
    and 8 words per bit array, a default array size of 3 and of 1000, and a largest array size of 5
    and of $2^{21}$, with the constants of stage~C1 changed too; so no proof depends on their
    values;
  \item every proof of the stage takes at most 20 lines, with named loop invariants and lemmas in
    the shape of the verification conditions, and \verb|Ratpoints/Proof/Words.lean| compiles in a
    few seconds.
\end{itemize}
